\documentclass[11pt]{article}
\usepackage[margin=1in]{geometry}
\usepackage{enumitem}
% =============================================================================
% Preambles/header.tex — shared LaTeX/Beamer preamble for this project
%
% Use from a Beamer lecture by prepending:
%     \input{header}
% and compiling with TEXINPUTS=../Preambles:$TEXINPUTS (the /compile-latex
% skill does this automatically).
%
% PALETTE CONTRACT. Color names defined here MUST match the names in
% Quarto/theme-template.scss so Beamer and Quarto renderings use the same
% palette. The scripts/check-palette-sync.sh script enforces this.
%
% To customize: edit the HEX values below AND the matching $variable in
% Quarto/theme-template.scss, then run scripts/check-palette-sync.sh.
% =============================================================================

% -----------------------------------------------------------------------------
% Engine detection + encoding
%
% Our /compile-latex skill uses XeLaTeX, where inputenc is unnecessary and
% can produce encoding warnings. Only load inputenc under pdfTeX.
% -----------------------------------------------------------------------------
\usepackage{iftex}
\ifPDFTeX
  \usepackage[utf8]{inputenc}
\fi

\usepackage{amsmath, amssymb, amsthm}
\usepackage{booktabs}
\usepackage{graphicx}

% -----------------------------------------------------------------------------
% PALETTE — mirrors Quarto/theme-template.scss variable names.
% Load xcolor and define colors BEFORE anything that references them
% (notably hyperref, hypersetup, and the Beamer color assignments).
% -----------------------------------------------------------------------------
\usepackage{xcolor}

\definecolor{primary-blue}{HTML}{012169}      % headings, accents
\definecolor{primary-gold}{HTML}{B9975B}      % emphasis, borders
\definecolor{highlight-yellow}{HTML}{F2A900}  % markers, alerts
\definecolor{light-bg}{HTML}{E8EDF5}          % title / section backgrounds
\definecolor{jet}{HTML}{1A1A1A}               % body text

% Semantic colors (used in diagrams and annotations)
\definecolor{positive}{HTML}{15803D}          % good / observed / identified
\definecolor{negative}{HTML}{B91C1C}          % bad / problematic / bias
\definecolor{neutral}{HTML}{525252}           % reference / context

% Highlight accents (match .hi-* classes in the SCSS)
\definecolor{hi-slate}{HTML}{314F4F}
\definecolor{hi-green}{HTML}{2E8B57}
\definecolor{hi-red}{HTML}{C41E3A}

% Semantic aliases so skill templates and snippets can write
% \textcolor{accent}{...} without hard-coding "primary-blue".
\colorlet{accent}{primary-blue}
\colorlet{accent2}{primary-gold}

% -----------------------------------------------------------------------------
% Hyperref — loaded AFTER xcolor + palette so link colors resolve.
% -----------------------------------------------------------------------------
\usepackage{hyperref}
\hypersetup{colorlinks=true, linkcolor=primary-blue, urlcolor=primary-blue,
            citecolor=primary-blue, pdfborder={0 0 0}}

% -----------------------------------------------------------------------------
% Course code — set once per deck (after \input{header}, before \begin{document})
% via \coursecode{CS401}. Shown in the footer on every slide so a deck's course
% is identifiable without opening the title page. Empty by default.
% -----------------------------------------------------------------------------
\makeatletter
\newcommand{\coursecode}[1]{\gdef\@coursecodetext{#1}}
\gdef\@coursecodetext{}
\makeatother

% -----------------------------------------------------------------------------
% Beamer theme assignments (only apply under Beamer).
%
% We detect Beamer via \@ifclassloaded{beamer}, which is the documented,
% non-internal way. This must be wrapped in \makeatletter/\makeatother
% because \@ifclassloaded contains an @-catcode character.
% -----------------------------------------------------------------------------
\makeatletter
\@ifclassloaded{beamer}{%
  \usefonttheme{professionalfonts}
  \setbeamercolor{structure}{fg=primary-blue}
  \setbeamercolor{frametitle}{fg=primary-blue}
  \setbeamercolor{title}{fg=primary-blue}
  \setbeamercolor{subtitle}{fg=primary-gold}
  \setbeamercolor{author}{fg=primary-blue}
  \setbeamercolor{institute}{fg=neutral}
  \setbeamercolor{date}{fg=primary-gold}
  \setbeamercolor{itemize item}{fg=primary-blue}
  \setbeamercolor{itemize subitem}{fg=primary-gold}
  \setbeamercolor{alerted text}{fg=primary-gold}
  \setbeamercolor{block title}{fg=white, bg=primary-blue}
  \setbeamercolor{block body}{bg=light-bg}
  % Semantic block variants: block=definition/concept (blue), exampleblock=
  % worked example (green/positive), alertblock=key takeaway or caution (gold).
  % Keep to <=2 colored boxes per slide across ALL three variants (INV-7).
  \setbeamercolor{block title example}{fg=white, bg=positive}
  \setbeamercolor{block body example}{bg=light-bg}
  \setbeamercolor{block title alerted}{fg=white, bg=primary-gold}
  \setbeamercolor{block body alerted}{bg=light-bg}

  % Minimal footer: course code (left) + page X / N (right), no navigation clutter
  \setbeamertemplate{navigation symbols}{}
  \setbeamertemplate{footline}{%
    \color{neutral}\footnotesize\hspace{0.7em}\@coursecodetext\hfill
    \insertframenumber/\inserttotalframenumber\hspace{0.7em}\vspace{0.3em}}
}{}
\makeatother

% -----------------------------------------------------------------------------
% TikZ setup — libraries the snippet gallery uses
% -----------------------------------------------------------------------------
\usepackage{tikz}
\usetikzlibrary{arrows.meta, positioning, calc, decorations.pathreplacing,
                fit, shapes.geometric, backgrounds}

% Shared TikZ styles — snippets and hand-written diagrams can pick these up
% via `\begin{tikzpicture}[every node/.style={...}]` or by naming specific
% styles (e.g., `\node[dag-node] (X) at (0,0) {X};`).
\tikzset{
  % Node styles for causal DAGs and similar
  dag-node/.style = {
    draw=primary-blue, fill=light-bg, rounded corners=2pt,
    minimum width=1.2cm, minimum height=0.9cm, align=center, font=\small
  },
  decision-node/.style = {
    draw=primary-gold, fill=white, diamond, aspect=1.8,
    minimum width=1.8cm, minimum height=1.2cm, align=center, font=\small,
    inner sep=1pt
  },
  % Edge styles — observed vs counterfactual
  observed-edge/.style = {-{Latex[length=2.5mm]}, thick, draw=jet},
  counterfactual-edge/.style = {-{Latex[length=2.5mm]}, thick, draw=neutral, dashed},
  confound-edge/.style = {-{Latex[length=2.5mm]}, thick, draw=negative, dashed},
  % Data-point styles for plots
  observed-dot/.style = {fill=primary-blue, draw=primary-blue, circle, inner sep=1.2pt},
  counterfactual-dot/.style = {fill=white, draw=neutral, circle, inner sep=1.2pt}
}

% -----------------------------------------------------------------------------
% Convenience macros
% -----------------------------------------------------------------------------
\newcommand{\muted}[1]{\textcolor{neutral}{#1}}
\newcommand{\key}[1]{\textcolor{primary-gold}{\textbf{#1}}}
\newcommand{\good}[1]{\textcolor{positive}{#1}}
\newcommand{\bad}[1]{\textcolor{negative}{#1}}

% Standout frame macro: full-bleed dark background for section transitions.
% Beamer-only (uses \begin{frame}/\setbeamercolor/\usebeamerfont) -- do not
% call from an article-class file (e.g. Notes/<CODE>/*-notes.tex). \muted,
% \key, \good, \bad, and \coursecode above are plain xcolor/gdef and are
% safe to use from either class; verified when Notes/article-class support
% was added.
% Expand: \transitionslide{Your transition title}
\newcommand{\transitionslide}[1]{%
  \begingroup
  \setbeamercolor{background canvas}{bg=primary-blue}
  \begin{frame}[plain]
    \centering
    \vfill
    {\usebeamerfont{title}\color{white}\Large #1\par}
    \vfill
  \end{frame}
  \endgroup
}

% Section divider frame (Beamer-only), an alternative to \transitionslide with a
% darker two-line style: near-black (jet) background, a small white label line
% above a larger gold title, and a thin gold rule along the bottom edge.
% Expand: \sectiondivider{Section 2}{Pointers}
\newcommand{\sectiondivider}[2]{%
  \begingroup
  \setbeamercolor{background canvas}{bg=jet}
  \begin{frame}[plain]
    \vspace*{3.0cm}
    {\color{white}\ttfamily\large #1\par}
    \vspace{0.5cm}
    {\color{primary-gold}\LARGE #2\par}
    \begin{tikzpicture}[remember picture, overlay]
      \draw[primary-gold, line width=2.5pt]
        ($(current page.south west)+(0,0.1)$) -- ($(current page.south east)+(0,0.1)$);
    \end{tikzpicture}
  \end{frame}
  \endgroup
}

% -----------------------------------------------------------------------------
% Channel watermark — "Concepts That Click" (YouTube).
%
% Article-class files (CompetitiveExam Q/A sets, Notes, Assignments) get a
% light diagonal draft watermark on every page plus a centered footer naming
% the channel. Beamer decks get a subtle TikZ background watermark instead
% (the footline already carries the course code). Centralized here so every
% MCQ PDF is branded without per-file edits.
% -----------------------------------------------------------------------------
\makeatletter
\@ifclassloaded{article}{%
  \usepackage{draftwatermark}
  \SetWatermarkText{Concepts That Click}
  \SetWatermarkScale{1}
  \SetWatermarkFontSize{2cm}
  \SetWatermarkLightness{0.86}
  \SetWatermarkAngle{45}
  \usepackage{fancyhdr}
  \pagestyle{fancy}
  \fancyhf{}
  \fancyfoot[C]{\footnotesize\textcolor{neutral}{Concepts That Click~|~YouTube: \href{https://www.youtube.com/@concepts.thatclick}{@concepts.thatclick}}}
  \renewcommand{\headrulewidth}{0pt}
  \renewcommand{\footrulewidth}{0pt}
  \fancypagestyle{plain}{%
    \fancyhf{}%
    \fancyfoot[C]{\footnotesize\textcolor{neutral}{Concepts That Click~|~YouTube: \href{https://www.youtube.com/@concepts.thatclick}{@concepts.thatclick}}}%
    \renewcommand{\headrulewidth}{0pt}%
    \renewcommand{\footrulewidth}{0pt}%
  }
}{}
\@ifclassloaded{beamer}{%
  \setbeamertemplate{background}{%
    \begin{tikzpicture}[remember picture, overlay]
      \node[rotate=45, opacity=0.055, align=center]
        at (current page.center)
        {\Huge\textcolor{neutral}{Concepts That Click}};
    \end{tikzpicture}%
  }
}{}
\makeatother 
\coursecode{JKSSB}
\renewcommand{\thesection}{60.\arabic{section}}
\newtheorem{example}{Example}[section]
\newenvironment{definitionbox}[1]{%
  \par\noindent\textbf{Definition (#1).}\ \itshape
}{\par\vspace{0.3em}}
\title{Current OS, Processes, Linux and Utilities --- Detailed Lecture Notes}
\author{JKSSB Computer Awareness}
\date{\today}

\begin{document}
\maketitle

\section{Deadlock and its four necessary conditions}
A \key{deadlock} is a situation in which a group of processes are each waiting
for a resource that is held by another process in the group, so that none of
them can ever proceed. A simple picture is two processes: process A holds a
printer and waits for a scanner, while process B holds the scanner and waits
for the printer. Each waits for the other, and both stay stuck forever.

For a deadlock to be possible, four conditions must hold \emph{at the same
time}. They are usually listed as follows.

\begin{center}
\begin{tabular}{p{0.28\textwidth}p{0.62\textwidth}}
\toprule
\textbf{Condition} & \textbf{Meaning} \\
\midrule
Mutual exclusion & A resource is held by only one process at a time. \\
Hold and wait & A process holds at least one resource while waiting for more. \\
No preemption & A resource cannot be forcibly taken away from its holder. \\
Circular wait & A cycle of processes exists, each waiting for a resource held by the next. \\
\bottomrule
\end{tabular}
\end{center}

\begin{definitionbox}{Deadlock}
A state in which a set of processes are each blocked waiting for a resource
held by another process in the same set, so no process can make progress.
\end{definitionbox}

\begin{definitionbox}{No preemption (as a condition)}
The rule that a resource already allocated to a process cannot be taken away
by force; it must be released voluntarily. Its opposite, \key{preemption}, is
used to \emph{recover} from a deadlock.
\end{definitionbox}

\section{Deadlock recovery by resource preemption}
Once a deadlock has been detected, the operating system must break it. One
recovery method is \key{resource preemption}: the system takes a resource away
from one of the processes and gives it to another process that needs it,
breaking the circular wait. The chosen process is usually rolled back to a
safe earlier state, or restarted. Preemption is also used to break the
\emph{no-preemption} condition directly, since it removes the guarantee that a
resource stays with its holder.

\begin{example}[The exam contrast]
An item asks for a method of \emph{recovering} from deadlock. The correct
answer is \good{resource preemption} --- the system forcibly reclaims a
resource. A distractor may name ``mutual exclusion'' or ``circular wait'',
but those are \bad{necessary conditions} for a deadlock, not recovery
methods. The other trap is to think preemption is one of the four conditions;
the condition is actually \key{no preemption}. (PYQ-10)
\end{example}

\section{Processes and the scheduler}
A \key{process} is a program in execution. When a program is launched, the
operating system creates a process for it and keeps the information needed to
run it --- its state, its program counter, and the resources it holds. Because
many processes are usually ready at once but only one can run on a given CPU
core at a time, the operating system includes a \key{scheduler} that decides
which ready process gets the CPU next.

Scheduling may be \key{non-preemptive}, where a running process keeps the CPU
until it finishes or blocks, or \key{preemptive}, where the operating system
can interrupt a running process and switch to another. Priorities are one of
the most common ways to choose between processes, and it is the combination of
\emph{priorities} and \emph{preemption} that produces the problems studied in
the next sections: priority inversion and starvation.

\begin{definitionbox}{Process}
A program in execution, together with the state and resources the operating
system maintains for it.
\end{definitionbox}

\begin{definitionbox}{Scheduler}
The operating system component that selects which ready process is given the
CPU next.
\end{definitionbox}

\section{Priority inversion versus priority inheritance}
In a priority-based system, a task with a higher priority is supposed to run
before one with a lower priority. \key{Priority inversion} breaks that rule.
It happens when a \key{high-priority} task is blocked because a
\key{low-priority} task currently holds a resource the high-priority task
needs. While the low-priority task holds that resource, a
\key{medium-priority} task can preempt it and run. The result is that the
high-priority task ends up waiting behind the medium-priority task, even
though it should have run first --- its priority has effectively been
inverted.

The standard fix is \key{priority inheritance}. While the low-priority task
holds the resource, it temporarily \key{inherits} the high priority of the
task that is waiting for it, so it runs at that high priority, releases the
resource quickly, and then returns to its original low priority. This lets the
high-priority task obtain the resource without waiting behind the
medium-priority task.

\begin{center}
\begin{tabular}{p{0.24\textwidth}p{0.32\textwidth}p{0.32\textwidth}}
\toprule
\textbf{Point} & \textbf{Priority inversion} & \textbf{Priority inheritance} \\
\midrule
What it is & A problem in which priority order is reversed. & A solution that temporarily raises the holder's priority. \\
Cause & A low-priority task holds a needed resource. & Applied while a resource is held and awaited. \\
Effect & High-priority task waits behind a medium-priority task. & Low-priority holder finishes fast, so the high-priority task proceeds. \\
\bottomrule
\end{tabular}
\end{center}

\begin{definitionbox}{Priority inversion}
A condition in which a high-priority task is indirectly preempted by a
lower-priority task, typically because a low-priority task holds a shared
resource.
\end{definitionbox}

\begin{example}[What does \emph{not} fix inversion]
An item asks which method prevents priority inversion. \key{Round-robin}
scheduling gives every ready process a turn but knows nothing about resource
holding, so by itself it does \bad{not} prevent inversion. The correct method
is \good{priority inheritance} (or the related priority-ceiling protocol).
(PYQ-11)
\end{example}

\section{Starvation and aging}
\key{Starvation} (indefinite blocking) is different from deadlock. A starved
process is \emph{ready} to run and could run in principle, but it is
\key{repeatedly passed over} by the scheduler, so it may wait for an
unbounded time. A common case is a low-priority process in a strict priority
system: as long as higher-priority processes keep arriving, the low-priority
process never gets the CPU.

The standard remedy for starvation is \key{aging}: the system gradually
\key{raises the priority} of a process the longer it waits, so that eventually
it becomes the highest-priority ready process and runs.

\begin{definitionbox}{Starvation}
A scheduling condition in which a ready process is indefinitely denied the
resources it needs because other processes are repeatedly preferred.
\end{definitionbox}

\begin{center}
\begin{tabular}{p{0.22\textwidth}p{0.34\textwidth}p{0.34\textwidth}}
\toprule
\textbf{Point} & \textbf{Starvation} & \textbf{Deadlock} \\
\midrule
Can the process ever run? & Yes, in principle & No, not without outside action \\
Cause & Scheduling choices (long wait) & A circular resource dependency \\
Cure & Aging (raise priority with waiting time) & Detection + recovery (for example preemption) \\
\bottomrule
\end{tabular}
\end{center}

\section{Preemptive priority scheduling}
In \key{priority scheduling}, the scheduler always picks the ready process with
the \key{highest priority}. The scheme is \key{preemptive} when a process that
is currently running can be \key{interrupted} as soon as a process with a
higher priority becomes ready; the higher-priority process then runs. It is
\key{non-preemptive} when the running process keeps the CPU until it either
finishes or blocks for input/output. Preemptive priority scheduling gives
quick service to important tasks, but with strict priorities it can starve
low-priority tasks unless aging is added.

\section{Real-time operating systems (RTOS)}
An \key{RTOS} is a \key{Real-Time Operating System}. It is designed so that
the system responds within a \key{bounded time}, allowing \key{deadlines} to be
met. Correctness in a real-time system depends not only on the \emph{result}
but also on \emph{when} the result is produced.

\begin{definitionbox}{Real-time operating system (RTOS)}
An operating system designed to process data and events within guaranteed
time limits, so that deadline requirements are satisfied.
\end{definitionbox}

Real-time systems are classified by what happens when a deadline is missed. In
a \key{hard real-time} system a missed deadline is a \bad{failure} and can have
serious or catastrophic consequences. In a \key{soft real-time} system a
missed deadline only \good{degrades} the quality of the result; the result is
still useful. The trap is to think an RTOS is simply ``a very fast operating
system''. Speed helps, but the defining property is \key{meeting deadlines}.

\section{The Linux kernel}
The \key{kernel} is the \key{core of the operating system}. It is the part
that runs with full control of the hardware and provides the basic services
every other program relies on. In Linux, the kernel manages
\key{processes} (through the scheduler), \key{memory}, \key{devices} (through
device drivers), and \key{file systems}.

The kernel runs in a privileged mode called \key{kernel space}, while ordinary
application programs run in \key{user space}. A user program asks the kernel
for a service --- reading a file, for example --- through a \key{system call};
the kernel performs the operation and returns the result. The Linux kernel is
\key{open-source}: its source code can be viewed, modified, and shared.

\begin{definitionbox}{Kernel}
The central component of an operating system that manages the hardware
resources (processes, memory, devices, and file systems) and mediates between
applications and the hardware.
\end{definitionbox}

The \key{shell} is frequently confused with the kernel. The shell is the
\emph{command interpreter} that a user types commands into; it is not the
kernel. The shell passes the user's requests to the kernel, which does the
actual work.

\section{\texttt{/root} versus \texttt{/home}}
Linux stores the personal files of each user in that user's \key{home
directory}. The \key{superuser}, whose account name is \key{root}, has its
home directory at \key{\texttt{/root}}. Ordinary users' home directories are
normally placed under \key{\texttt{/home}}, for example
\texttt{/home/alice}.

It is important not to confuse two different things with similar names:

\begin{center}
\begin{tabular}{p{0.18\textwidth}p{0.70\textwidth}}
\toprule
\textbf{Path} & \textbf{Meaning} \\
\midrule
\texttt{/} & The \key{root of the whole file system} --- the single top directory from which all other directories branch. \\
\texttt{/root} & The \key{home directory of the root (superuser)} account. \\
\texttt{/home} & The parent directory holding the home directories of \key{ordinary users}. \\
\bottomrule
\end{tabular}
\end{center}

The trap is to read ``\texttt{/root}'' as ``the root of the file system''. The
top of the file system is \texttt{/} (a single slash). \texttt{/root} is only
the superuser's own home folder.

\section{File systems}
A \key{file system} is the method an operating system uses to organise, store,
and retrieve files and folders on a storage device. Different systems use
different file systems.

\begin{center}
\begin{tabular}{p{0.16\textwidth}p{0.22\textwidth}p{0.50\textwidth}}
\toprule
\textbf{File system} & \textbf{Full form} & \textbf{Typical use} \\
\midrule
NTFS & New Technology File System & Default for modern Windows installations. \\
FAT32 & File Allocation Table, 32-bit & Older system; widely compatible with small drives and memory cards. \\
ext4 & Fourth Extended file system & Native and common default file system for Linux. \\
\bottomrule
\end{tabular}
\end{center}

The exam cue is the pairing: \key{NTFS = Windows}, \key{ext4 = Linux}, and
\key{FAT32 = broad compatibility}. Full forms are commonly asked directly.

\section{Windows utilities: Store, WordPad/Notepad, Disk Cleanup}
The \key{Windows Store} (now the \key{Microsoft Store}) is the built-in
\key{online marketplace} in Windows. Users \key{search, download, install, and
update} apps from it, so applications come from a trusted catalogue rather
than from arbitrary websites.

Windows includes two simple text editors that are easy to confuse.

\begin{center}
\begin{tabular}{p{0.14\textwidth}p{0.36\textwidth}p{0.36\textwidth}}
\toprule
\textbf{Point} & \textbf{Notepad} & \textbf{WordPad} \\
\midrule
Text handling & \key{Plain text} only; no formatting & \key{Basic formatting}: bold, italic, fonts, colour \\
Typical save & \texttt{.txt} & \texttt{.rtf} (rich text format) \\
Capability & The smallest, fastest text editor & Can insert pictures and objects \\
\bottomrule
\end{tabular}
\end{center}

\key{Notepad} is a plain-text editor: it displays and saves characters without
any formatting. \key{WordPad} supports \emph{basic} rich-text formatting and
can include pictures, but it is still lighter than a full word processor such
as MS Word.

\key{Disk Cleanup} is a built-in Windows utility that \key{frees disk space}.
It deletes files that are not needed, such as temporary files, thumbnail
caches, downloaded program files, and the contents of the \key{Recycle Bin}.
It is the tool to reach for when a drive is running short of free space.

\section{Defragmenter and the SSD caution}
On a traditional \key{hard disk drive (HDD)}, a file can become
\key{fragmented}: its pieces are scattered in different places on the rotating
disk, which makes reading slower. The \key{Disk Defragmenter} rearranges the
pieces so that each file is stored \key{contiguously} (in one run), which
improves access speed.

\textbf{SSD caution.} A \key{solid-state drive (SSD)} has \key{no moving
parts}. Defragmenting it gives no speed benefit and only adds
\key{unnecessary writes}, which reduce the drive's working life. SSDs should
\bad{not} be defragmented; Windows uses the \key{TRIM} command for SSDs
instead.

\begin{definitionbox}{Defragmentation}
The process of rearranging the fragments of files on a hard disk so that each
file occupies contiguous space, improving read performance.
\end{definitionbox}

\section{Imaging Fax and the \texttt{xcopy /s} command}
\key{Windows Fax and Scan} is the built-in \key{imaging and fax} utility in
Windows. With a scanner it can \key{scan} a document and save it as an image;
with a fax modem or fax service it can \key{send and receive faxes}. This is
the tool the exam means by ``Imaging Fax''.

\key{xcopy} is a Windows \key{command-line} command that copies files and
whole directory trees. Its switch \key{\texttt{/s}} copies the named directory
\emph{including its subdirectories}, but it does not copy empty ones. Adding
\key{\texttt{/e}} copies empty subdirectories as well. The exact form to
remember is \key{\texttt{xcopy /s}}: ``copy with subdirectories''.

\begin{center}
\begin{tabular}{p{0.20\textwidth}p{0.68\textwidth}}
\toprule
\textbf{Command / tool} & \textbf{What it does} \\
\midrule
\texttt{xcopy /s} & Copies files and subdirectories (empty ones excluded). \\
\texttt{xcopy /s /e} & Also copies empty subdirectories. \\
Windows Fax and Scan & Scans documents and sends/receives faxes (Imaging Fax). \\
\bottomrule
\end{tabular}
\end{center}

\section{Exam retrieval and common confusions}
Six distinctions prevent most errors on this topic.
\begin{enumerate}
  \item Deadlock's four conditions are mutual exclusion, hold and wait,
    \key{no preemption}, and circular wait. \key{Preemption} is a
    \key{recovery} method, not a condition (PYQ-10).
  \item \key{Priority inversion} is the problem; \key{priority inheritance} is
    the cure. Round-robin by itself does \bad{not} prevent inversion (PYQ-11).
  \item \key{Starvation} is indefinite \emph{scheduling} delay, not deadlock;
    \key{aging} fixes starvation.
  \item An \key{RTOS} is defined by \key{meeting deadlines} (hard vs soft), not
    merely by being fast.
  \item The Linux \key{kernel} is the core; the \key{shell} is the command
    interpreter. \texttt{/root} is the superuser's home; \texttt{/} is the top
    of the file system; \texttt{/home} holds ordinary users' homes.
  \item \key{Notepad} = plain text; \key{WordPad} = basic rich text. Disk
    Cleanup frees space; \key{never defragment an SSD} (use TRIM);
    \texttt{xcopy /s} copies subdirectories.
\end{enumerate}

\begin{example}[Tracing a mixed item]
An item lists statements and asks which are true:
(I) Preemption is one of the four necessary conditions for deadlock.
(II) Priority inheritance reduces priority inversion.
(III) Aging helps prevent starvation.
Statement I is \bad{false}: preemption is a recovery method; the condition is
\emph{no} preemption. Statements II and III are \good{true}. The answer is
``II and III only''.
\end{example}

\subsection*{Exercises}
\begin{enumerate}[label=\arabic*.]
  \item List the four necessary conditions for deadlock, and name one method
    of recovering from a deadlock.
  \item Explain priority inversion and state which technique resolves it.
  \item Distinguish starvation from deadlock, and name the standard remedy for
    starvation.
  \item What is an RTOS? Give the difference between a hard and a soft
    real-time system.
  \item State the meaning of \texttt{/root} and \texttt{/home} in Linux, and
    name the default file system of Windows and of Linux.
  \item Explain why an SSD should not be defragmented, and state what
    \texttt{xcopy /s} does.
\end{enumerate}

\subsection*{Solutions}
\begingroup\small
\begin{enumerate}[label=\arabic*.]
  \item Mutual exclusion, hold and wait, no preemption, and circular wait. A
    deadlock can be recovered from by resource preemption (taking a resource
    away from one process and giving it to another).
  \item Priority inversion is when a high-priority task waits because a
    low-priority task holds a needed resource, while a medium-priority task
    runs. \key{Priority inheritance} resolves it: the low-priority holder
    temporarily inherits the high priority and releases the resource quickly.
  \item A starved process is ready but repeatedly passed over, so it may wait
    indefinitely; a deadlocked process can never proceed because of a circular
    resource dependency. \key{Aging} --- gradually raising the waiting
    process's priority --- prevents starvation.
  \item An RTOS (Real-Time Operating System) guarantees a response within a
    bounded time so deadlines are met. In a \key{hard} real-time system a
    missed deadline is a failure; in a \key{soft} real-time system a missed
    deadline only degrades quality.
  \item \texttt{/root} is the home directory of the superuser (root);
    \texttt{/home} holds the home directories of ordinary users. Windows uses
    \key{NTFS}; Linux uses \key{ext4}.
  \item An SSD has no moving parts, so defragmenting gives no speed benefit and
    only adds wear; SSDs use \key{TRIM} instead. \texttt{xcopy /s} copies a
    directory including its subdirectories, excluding empty ones.
\end{enumerate}
\endgroup
\end{document}
